% =============================================================================
\chapter{Dynamic Programming}
\label{ch:dynamic_programming}
% =============================================================================

The name ``dynamic programming'' is one of the great misnomers in mathematics. Richard Bellman, who invented the framework in the 1950s at the RAND Corporation, later confessed that he chose the word ``dynamic'' to impress his Pentagon sponsors and ``programming'' because it sounded good---the actual content has nothing to do with computer programming in the modern sense. What Bellman discovered was something far deeper: a \emph{principle of optimality} stating that an optimal policy has the property that, whatever the initial state and decision, the remaining decisions must constitute an optimal policy from the resulting state onward.

This recursive insight transforms an impossibly large optimisation problem---choosing the best action at every state for all time---into a sequence of local computations. When the environment is fully known (transition probabilities and rewards are given), dynamic programming solves Markov decision processes exactly, through two complementary algorithms: \emph{policy iteration} and \emph{value iteration}. These are not merely historical curiosities; they are the theoretical bedrock upon which all of reinforcement learning is built. Every modern RL algorithm, from Q-learning to policy gradient methods, can be understood as an approximation of the ideas in this chapter.

\begin{intuition}
Dynamic programming (DP) solves MDPs when the transition model $P(s'|s,a)$ and
reward function $R(s,a,s')$ are perfectly known. It constitutes the theoretical
foundation upon which all reinforcement learning methods rest.
\end{intuition}

% -----------------------------------------------------------------------------
\section{Principle of Dynamic Programming}
% -----------------------------------------------------------------------------

Dynamic programming, introduced by Richard Bellman in the 1950s, exploits the
recursive structure of the Bellman equations to iteratively compute value
functions. The two fundamental algorithms are:
\begin{enumerate}
    \item \textbf{Policy evaluation}: computing $V^\pi$ for a given policy $\pi$.
    \item \textbf{Policy improvement}: constructing a better policy from $V^\pi$.
\end{enumerate}

% -----------------------------------------------------------------------------
\section{Policy Evaluation}
% -----------------------------------------------------------------------------

\begin{definition}[Iterative Policy Evaluation]
Iterative policy evaluation computes $V^\pi$ by repeatedly applying the
Bellman operator $\mathcal{T}^\pi$:
\[
    V_{k+1}(s) = \sum_{a} \pi(a|s) \sum_{s'} P(s'|s,a)[R(s,a,s') + \gamma V_k(s')]
    \quad \forall s \in \mathcal{S}.
\]
\end{definition}

\begin{theorem}[Convergence of Policy Evaluation]
\label{thm:conv_eval}
For any initialization $V_0$, the sequence $(V_k)_{k \geq 0}$ defined by
$V_{k+1} = \mathcal{T}^\pi V_k$ converges to $V^\pi$ as $k \to \infty$.
Moreover:
\[
    \norm{V_k - V^\pi}_\infty \leq \gamma^k \norm{V_0 - V^\pi}_\infty.
\]
\end{theorem}

\begin{proof}
This is a direct consequence of Theorem~\ref{thm:contraction}. The operator
$\mathcal{T}^\pi$ is a $\gamma$-contraction in the infinity norm, so by
Banach's fixed-point theorem, the sequence converges geometrically to the
unique fixed point $V^\pi$.
\end{proof}

\begin{algorithme}[title=\textbf{Iterative Policy Evaluation}]
\begin{enumerate}
    \item Initialize $V(s) \leftarrow 0$ for all $s \in \mathcal{S}$
    \item \textbf{Repeat} until convergence ($\Delta < \theta$):
    \begin{enumerate}
        \item $\Delta \leftarrow 0$
        \item For each $s \in \mathcal{S}$:
        \begin{enumerate}
            \item $v \leftarrow V(s)$
            \item $V(s) \leftarrow \sum_a \pi(a|s) \sum_{s'} P(s'|s,a)[R(s,a,s') + \gamma V(s')]$
            \item $\Delta \leftarrow \max(\Delta, |v - V(s)|)$
        \end{enumerate}
    \end{enumerate}
    \item Return $V \approx V^\pi$
\end{enumerate}
\end{algorithme}

% -----------------------------------------------------------------------------
\section{Policy Improvement}
% -----------------------------------------------------------------------------

\begin{theorem}[Policy Improvement Theorem]
\label{thm:policy_improvement}
Let $\pi$ be a policy and $\pi'$ the greedy policy with respect to $V^\pi$:
\[
    \pi'(s) = \argmax_{a \in \mathcal{A}} \sum_{s'} P(s'|s,a)[R(s,a,s') + \gamma V^\pi(s')].
\]
Then $V^{\pi'}(s) \geq V^\pi(s)$ for all $s \in \mathcal{S}$, with equality
if and only if $\pi$ is already optimal.
\end{theorem}

\begin{proof}
For any state $s$:
\begin{align*}
    V^\pi(s) &\leq \max_a Q^\pi(s,a) = Q^\pi(s, \pi'(s)) \\
    &= \sum_{s'} P(s'|s,\pi'(s))[R(s,\pi'(s),s') + \gamma V^\pi(s')] \\
    &\leq \sum_{s'} P(s'|s,\pi'(s))[R(s,\pi'(s),s') + \gamma Q^\pi(s',\pi'(s'))] \\
    &\leq \cdots \leq V^{\pi'}(s).
\end{align*}
The result follows by recursive application of the inequality.
\end{proof}

% -----------------------------------------------------------------------------
\section{Policy Iteration}
% -----------------------------------------------------------------------------

\begin{algorithme}[title=\textbf{Policy Iteration}]
\begin{enumerate}
    \item Initialize $\pi(s)$ arbitrarily for all $s$
    \item \textbf{Repeat}:
    \begin{enumerate}
        \item \textbf{Evaluation}: compute $V^\pi$ via iterative policy evaluation
        \item \textbf{Improvement}: for each $s$:
        \[
            \pi'(s) \leftarrow \argmax_a \sum_{s'} P(s'|s,a)[R(s,a,s') + \gamma V^\pi(s')]
        \]
        \item If $\pi' = \pi$, \textbf{stop} (optimal policy found)
        \item $\pi \leftarrow \pi'$
    \end{enumerate}
    \item Return $\pi^* = \pi$
\end{enumerate}
\end{algorithme}

\begin{theorem}[Convergence of Policy Iteration]
Policy iteration converges to the optimal policy $\pi^*$ in a finite number
of iterations (at most $|\mathcal{A}|^{|\mathcal{S}|}$ iterations, since
the number of deterministic policies is finite and each iteration strictly
improves the policy).
\end{theorem}

% -----------------------------------------------------------------------------
\section{Value Iteration}
% -----------------------------------------------------------------------------

\begin{definition}[Value Iteration]
Value iteration combines evaluation and improvement in a single update:
\[
    V_{k+1}(s) = \max_{a \in \mathcal{A}} \sum_{s'} P(s'|s,a)
    [R(s,a,s') + \gamma V_k(s')] \quad \forall s \in \mathcal{S}.
\]
\end{definition}

\begin{theorem}[Convergence of Value Iteration]
The sequence $(V_k)$ defined by value iteration converges to $V^*$:
\[
    \norm{V_k - V^*}_\infty \leq \gamma^k \norm{V_0 - V^*}_\infty.
\]
The number of iterations to reach precision $\epsilon$ is:
\[
    k \geq \frac{\log(\norm{V_0 - V^*}_\infty / \epsilon)}{\log(1/\gamma)}.
\]
\end{theorem}

\begin{algorithme}[title=\textbf{Value Iteration}]
\begin{enumerate}
    \item Initialize $V(s) \leftarrow 0$ for all $s \in \mathcal{S}$
    \item \textbf{Repeat} until convergence ($\Delta < \theta$):
    \begin{enumerate}
        \item $\Delta \leftarrow 0$
        \item For each $s \in \mathcal{S}$:
        \begin{enumerate}
            \item $v \leftarrow V(s)$
            \item $V(s) \leftarrow \max_a \sum_{s'} P(s'|s,a)[R(s,a,s') + \gamma V(s')]$
            \item $\Delta \leftarrow \max(\Delta, |v - V(s)|)$
        \end{enumerate}
    \end{enumerate}
    \item Extract policy: $\pi^*(s) = \argmax_a \sum_{s'} P(s'|s,a)[R + \gamma V(s')]$
    \item Return $V^*, \pi^*$
\end{enumerate}
\end{algorithme}

% -----------------------------------------------------------------------------
\section{Comparison of Methods}
% -----------------------------------------------------------------------------

\begin{center}
\renewcommand{\arraystretch}{1.3}
\begin{tabular}{lcc}
    \toprule
    & \textbf{Policy Iteration} & \textbf{Value Iteration} \\
    \midrule
    Cost per iteration & $O(|\mathcal{S}|^3 + |\mathcal{S}|^2 |\mathcal{A}|)$ &
    $O(|\mathcal{S}|^2 |\mathcal{A}|)$ \\
    Number of iterations & Few (often $< 10$) & More ($O(1/\log(1/\gamma))$) \\
    Memory & $V + \pi$ & $V$ only \\
    Convergence & Finite, exact & Asymptotic \\
    \bottomrule
\end{tabular}
\end{center}

\begin{remark}
In practice, policy iteration often converges in very few iterations (5--10),
even for large MDPs. Value iteration requires more iterations but each is
less expensive since it avoids the full policy evaluation step.
\end{remark}

% -----------------------------------------------------------------------------
\section{Generalized Policy Iteration (GPI)}
% -----------------------------------------------------------------------------

\begin{definition}[GPI --- Generalized Policy Iteration]
\textbf{Generalized Policy Iteration} refers to any alternation between
(partial or complete) policy evaluation and policy improvement. Nearly all
RL algorithms can be viewed as instances of GPI.
\end{definition}

\begin{remark}
Value iteration is a limiting case of GPI where evaluation performs only a
single Bellman backup before improvement.
\end{remark}

% -----------------------------------------------------------------------------
\section{Python Implementation}
% -----------------------------------------------------------------------------

\begin{codeblock}[title=\textbf{Policy Evaluation}]
\begin{minted}{python}
import numpy as np

def policy_evaluation(P, R, pi, gamma=0.99, theta=1e-8):
    """
    Iterative policy evaluation.
    P: transition matrix (nS, nA, nS)
    R: reward matrix (nS, nA, nS)
    pi: deterministic policy (nS,) -> action indices
    """
    nS = P.shape[0]
    V = np.zeros(nS)

    while True:
        delta = 0.0
        for s in range(nS):
            a = pi[s]
            v = V[s]
            V[s] = np.sum(P[s, a, :] * (R[s, a, :] + gamma * V))
            delta = max(delta, abs(v - V[s]))
        if delta < theta:
            break
    return V
\end{minted}
\end{codeblock}

\begin{codeblock}[title=\textbf{Policy Iteration}]
\begin{minted}{python}
def policy_iteration(P, R, gamma=0.99, theta=1e-8):
    """Complete policy iteration."""
    nS, nA = P.shape[0], P.shape[1]
    pi = np.zeros(nS, dtype=int)

    while True:
        V = policy_evaluation(P, R, pi, gamma, theta)
        stable = True
        for s in range(nS):
            old_action = pi[s]
            Q_s = np.array([
                np.sum(P[s, a, :] * (R[s, a, :] + gamma * V))
                for a in range(nA)
            ])
            pi[s] = np.argmax(Q_s)
            if old_action != pi[s]:
                stable = False
        if stable:
            break
    return V, pi
\end{minted}
\end{codeblock}

\begin{codeblock}[title=\textbf{Value Iteration}]
\begin{minted}{python}
def value_iteration(P, R, gamma=0.99, theta=1e-8):
    """Value iteration."""
    nS, nA = P.shape[0], P.shape[1]
    V = np.zeros(nS)

    while True:
        delta = 0.0
        for s in range(nS):
            v = V[s]
            Q_s = np.array([
                np.sum(P[s, a, :] * (R[s, a, :] + gamma * V))
                for a in range(nA)
            ])
            V[s] = np.max(Q_s)
            delta = max(delta, abs(v - V[s]))
        if delta < theta:
            break

    pi = np.zeros(nS, dtype=int)
    for s in range(nS):
        Q_s = np.array([
            np.sum(P[s, a, :] * (R[s, a, :] + gamma * V))
            for a in range(nA)
        ])
        pi[s] = np.argmax(Q_s)
    return V, pi
\end{minted}
\end{codeblock}

\begin{codeblock}[title=\textbf{Application to FrozenLake}]
\begin{minted}{python}
import gymnasium as gym

env = gym.make("FrozenLake-v1", is_slippery=True)
nS = env.observation_space.n  # 16
nA = env.action_space.n       # 4
P_mat = np.zeros((nS, nA, nS))
R_mat = np.zeros((nS, nA, nS))

for s in range(nS):
    for a in range(nA):
        for prob, next_s, reward, done in env.unwrapped.P[s][a]:
            P_mat[s, a, next_s] += prob
            R_mat[s, a, next_s] = reward

V_star, pi_star = value_iteration(P_mat, R_mat, gamma=0.99)
print("Optimal policy (4x4 grid):")
actions = ['<', 'v', '>', '^']
for i in range(4):
    print([actions[pi_star[4*i + j]] for j in range(4)])
\end{minted}
\end{codeblock}

% -----------------------------------------------------------------------------
\section{Exercises}
% -----------------------------------------------------------------------------

\begin{exercise}[Value Iteration Convergence]
Show that for the $4 \times 4$ gridworld from Chapter~\ref{ch:mdp} with
$\gamma = 0.9$ and $V_0 = 0$, value iteration achieves precision
$\epsilon = 10^{-6}$ in at most $k$ iterations. Compute $k$.
\end{exercise}

\begin{exercise}[Modified Policy Iteration]
Implement a version of policy iteration where evaluation performs exactly $m$
sweeps instead of converging. Study the effect of $m$ on overall convergence
speed for $m \in \{1, 3, 10, 100\}$.
\end{exercise}

\begin{exercise}[Asynchronous Sweeps]
Implement an asynchronous version of value iteration where states are updated
in random order (one state per iteration). Compare convergence with the
standard synchronous version.
\end{exercise}

\begin{exercise}[Matrix Solution]
For a 3-state MDP with the uniform policy, compute $V^\pi$ by matrix inversion
$(I - \gamma P^\pi)^{-1} r^\pi$ and verify that the result matches iterative
evaluation.
\end{exercise}
